
\subsubsection{5.1 Infrastructure Validation}

All 15 implementation tasks completed with 100\% Software Design Document (SDD) compliance in one Coder-Validator review cycle. All 67 unit and integration tests passed with zero failures. The prescreening pipeline ran end-to-end without errors. Table 1 summarizes the implementation status.

\textbf{Table 1. Implementation Completeness (h-e1)}

\begin{table}[htbp]
\centering
\resizebox{\columnwidth}{!}{%
\begin{tabular}{lll}
\toprule
Component & Status & Detail \\
\midrule
prescreening.py & Complete & 371 lines; main pipeline orchestration \\
evaluate.py & Complete & Per-problem evaluation and gate metrics \\
reward\_fn.py & Complete & R\_ratio, R\_binary, S\_term implementations \\
data\_loader.py & Complete & APPS loader (difficulty=0, $T\geq 3$ filter) \\
execution\_sandbox.py & Complete & Subprocess isolation, 5-second timeout \\
visualization.py & Complete & Gate metric and distribution plots \\
Tasks completed & 15/15 & 100\% SDD compliance \\
Tests passed & 67/67 & 0 failures \\
Coder-Validator cycles & 1 & No critical issues flagged \\
Problems loaded & 1,923 & Introductory split, $T\geq 3$ \\
Problems processed & 300 & seed=42, ~42 minutes on H100 NVL \\
\bottomrule
\end{tabular}
}
\end{table}

\subsubsection{5.2 Gate Metric Results}

Table 2 reports all Stage 1 gate metrics against pre-registered thresholds. All three quantitative gate criteria failed.

\textbf{Table 2. h-e1 Gate Metrics: Actual vs. Pre-registered Thresholds}

\begin{table}[htbp]
\centering
\resizebox{\columnwidth}{!}{%
\begin{tabular}{llll}
\toprule
Metric & Threshold & Actual & Pass/Fail \\
\midrule
fraction\_k\_pass\_ge1 & $\geq$ 0.10 & 0.0 & FAIL \\
pct\_groups\_above\_1.5x & $\geq$ 0.80 & 0.0 & FAIL \\
n\_prescreened & $\geq$ 50 & 0 & FAIL \\
Problems with S\_term $\in$ [0.3, 0.55] & $\geq$ 50 & 0/300 & FAIL \\
Implementation tasks & 15/15 & 15/15 & PASS \\
Integration tests & 67/67 & 67/67 & PASS \\
**Overall gate (MUST\_WORK)** & **PASS** & **PARTIAL** & --- \\
\bottomrule
\end{tabular}
}
\end{table}

The gate outcome is PARTIAL: infrastructure metrics pass; all quantitative behavioral metrics fail. The PARTIAL classification reflects the structural distinction between mechanism correctness (confirmed) and model behavioral prerequisite (unmet).

Figure 4 shows the gate metric comparison.

\begin{figure}[htbp]
\centering
\includegraphics[width=\columnwidth]{figures/fig_gate_metrics.png}
\caption{Figure 4: h-e1 gate metric results}
\label{fig:fig_gate_metrics}
\end{figure}

\textit{Figure 4. h-e1 gate metric evaluation. All three quantitative criteria (fraction\_k\_pass\_ge1, pct\_groups\_above\_1.5x, n\_prescreened) register at 0.0 against thresholds of 0.10, 0.80, and 50, respectively. Infrastructure metrics (tasks and tests) meet their targets exactly. The gap is attributable to the absent SFT checkpoint.}

\subsubsection{5.3 Base Model Pass Rate}

The central quantitative result is that Qwen2.5-Coder-7B-Instruct achieved 0\% pass rate on APPS introductory problems. Across 300 problems and k=8 rollouts per problem, 2,400 total solution attempts were executed. S\_term = 0.0 for all 300 problems: no rollout passed a single test case on any problem. The per-problem results CSV (h-e1/results/per\_problem\_results.csv) confirms this uniformly: every row records S\_term = 0.0 and tests\_passed\_vec = [0, 0, 0, 0, 0, 0, 0, 0] across all 8 rollouts.

This finding is notable because Qwen2.5-Coder-7B-Instruct achieves substantial performance on standard code generation benchmarks (approximately 72\% HumanEval pass@1). The APPS introductory tier is described as the easiest difficulty level in the benchmark. A probable contributing factor is format incompatibility: the instruction-tuned variant generates chat-format responses that include natural language explanation and markdown code fences, while the execution sandbox expects raw executable Python. The APPS execution harness does not strip markdown formatting prior to execution, so responses containing explanation text fail immediately at the Python parsing stage.

This interpretation is consistent with the finding of Du et al. [2025] that SFT initialization on APPS-format data is a prerequisite for non-trivial pass rates under execution-based evaluation. SFT training on APPS solutions teaches the model to produce bare executable Python in the expected format. The 0\% result does not distinguish between format incompatibility and genuine computational incapability; this distinction requires a targeted format diagnostic experiment (not conducted in this study).

\subsubsection{5.4 Simulated Advantage Distributions}

Because no empirical GRPO training data was generated (Stage 2 was blocked), Figure 3 shows simulated GRPO advantage distributions. The simulation uses 500 groups, q = 0.45, T = 5, G = 8, with rewards drawn from the Binomial model described in Section 3.1.

Under R\_binary, the advantage distribution is near-bimodal. With T = 5 test cases and per-test probability q = 0.45, the probability of passing all 5 tests is 0.$45^5 \approx$ 0.018. Most rollouts receive reward 0; a small fraction receives reward 1. The resulting advantage distribution clusters around two mass points.

Under R\_ratio, the advantage distribution is graded. Each rollout can pass 0 through 5 tests, yielding R\_ratio $\in$ {0.0, 0.2, 0.4, 0.6, 0.8, 1.0}. The distribution of advantages is more spread, with multiple distinct levels represented in most groups.

These simulations reflect the theoretical prediction and confirm the mathematical properties of the two reward functions under the Binomial model. They do not constitute empirical evidence from actual GRPO training, which remains to be conducted.

\begin{figure}[htbp]
\centering
\includegraphics[width=\columnwidth]{figures/fig_advantage_distribution.png}
\caption{Figure 3: Simulated GRPO advantage distributions}
\label{fig:fig_advantage_distribution}
\end{figure}

\textit{Figure 3. Simulated GRPO advantage distributions for R\_ratio (left) vs. R\_binary (right) under matched conditions: 500 groups, q = 0.45, T = 5, G = 8. These are simulation results, not empirical training observations. R\_ratio produces graded distributions with multiple distinct advantage levels; R\_binary produces a near-bimodal distribution concentrated at two values.}

\subsubsection{5.5 Stage 2 Status and Projected Outcomes}

Sub-hypotheses h-m1 through h-m4 are all blocked. No GRPO training was executed. The analytical projections for these hypotheses, derived from the Binomial variance model, are listed in Table 3 for reference. These projections are not experimental results; they are consequences of the mathematical analysis in Section 3.1 conditional on Stage 2 becoming executable.

\textbf{Table 3. Stage 2 Sub-hypotheses: Status and Analytical Projections}

\begin{table*}[htbp]
\centering
\resizebox{\dimexpr 2\columnwidth+\columnsep\relax}{!}{%
\begin{tabular}{lllll}
\toprule
Sub-hypothesis & Metric & Prediction & Analytical Basis & Experimental Status \\
\midrule
h-m1 & Distinct advantage levels per group & $\geq 5$ (R\_ratio) vs. ~2 (R\_binary) & T=5 yields 6 possible R\_ratio values & BLOCKED \\
h-m2 & $Cov(r_i, \|\nabla \theta log \pi (o_{i})\|)$ & Higher under R\_ratio & Higher reward variance $\rightarrow$ larger gradient steps & BLOCKED \\
h-m3 & Gradient SNR $\|E[A_{i}]\|/std(A_{i})$ in first 25\% training & $\geq 1.5\times$ under R\_ratio & Binomial variance ratio $\geq 1.5\times$ in window & BLOCKED \\
h-m4 & ZRF\_ratio(t*) vs. ZRF\_binary(t*) at t*=25\% training & ZRF\_ratio < 0.8 $\times$ ZRF\_binary; log-rank p<0.05 & Earlier advantage heterogeneity $\rightarrow$ earlier escape & BLOCKED \\
\bottomrule
\end{tabular}
}
\end{table*}

All Stage 2 metrics are pending availability of an SFT-initialized model checkpoint.

